\section{C语言程序的常见错误与调试}

\begin{frame}{HW1累加操作：函数里改了，外面为什么没变？}
    将求和操作放入一个C函数，用固定的小数据观察参数传递：
    \begin{Code}*{c}
        \lstinputlisting[firstline=4,lastline=6,firstnumber=4]{code/c/sum.c}
    \end{Code}
    \begin{Code}*{c}
        \lstinputlisting[firstline=8,lastline=13,firstnumber=8]{code/c/sum.c}
    \end{Code}
    \begin{itemize}
        \item 预期总和\texttt{8}，实际仍为\texttt{3}。
        \item 函数里的加法执行了；接下来应该观察哪一个\texttt{sum}？
    \end{itemize}
\end{frame}

\begin{frame}{HW1累加示例：为什么没改成功？怎样修复？}
    \begin{columns}[T]
        \begin{column}{0.48\textwidth}
            \small 本次实测将加数改为3：预期6，实际3。
            \lstset{basicstyle=\ttfamily\fontsize{6.6}{7.8}\selectfont,
                numbers=none,xleftmargin=0pt,breakatwhitespace=true}
            \begin{Code}*{plain}
                \lstinputlisting[linerange={1-10,13-16}]{code/output/sum-value-user-session.txt}
            \end{Code}
            \footnotesize \alert{原因：修改的是形参副本。}
            \par 函数内的\texttt{sum}变为6；\texttt{finish}回到\texttt{main}后，原变量仍为3。
        \end{column}
        \begin{column}{0.48\textwidth}
            \small 正确修改：传入地址，修改原变量。
            \lstset{basicstyle=\ttfamily\fontsize{7.4}{8.7}\selectfont,
                numbers=left,numberstyle=\fontsize{4.8}{5.5}\selectfont\color{gray},
                xleftmargin=6pt,breakatwhitespace=true}
            \begin{Code}*{c}
                \lstinputlisting[lastline=13]{code/fixed/sum-three.c}
            \end{Code}
            \footnotesize \texttt{\&sum}传入原变量的地址，\texttt{*sum}访问并修改它；输出\alert{6}。
            \par\medskip
            \scriptsize 完整修复版：\par\texttt{code/fixed/sum-three.c}。
        \end{column}
    \end{columns}
    \note{这是根据HW1累加操作构造的值传递课堂示例，不是原始作业错误程序。用户本次将前一页的加数5改为3，因此预期为6；保留用户实测GDB记录中的step、next、finish和3到6再回到3的观察值。修复版使用int *sum、调用处&sum、函数内*sum += x；指针本身仍按值传递，但通过复制的地址可以访问原变量。也可以将函数改为返回新总和，再由调用者赋值接收。}
\end{frame}

\begin{frame}{C的参数按值传递：怎样修改原变量？}
    \begin{center}
        \begin{tikzpicture}[>=Stealth,node distance=1.2cm,
            box/.style={draw=VSPPrimary,rounded corners,minimum width=3.5cm,
            minimum height=1cm,align=center,font=\small}]
            \node[box] (a) {调用者的\texttt{sum}\\仍为3};
            \node[box,right=of a] (b) {函数形参\texttt{sum}\\从3变为8};
            \draw[->,thick] (a.north east) to[bend left=15]
                node[above] {值传递：复制} (b.north west);
            \node[below=0.4cm of a,font=\small] {原变量没有改变};
            \node[below=0.4cm of b,font=\small] {修改的是局部副本};
        \end{tikzpicture}
    \end{center}
    \begin{itemize}
        \item GDB中在\texttt{addValue}内观察\texttt{sum}，再\texttt{finish}回到调用者比较。
        \item 需要修改原变量时，可传入地址\texttt{\&sum}，用形参\texttt{int *sum}接收，再通过\texttt{*sum}修改。
        \item 另一种设计是返回新总和，由调用者接收返回值。
    \end{itemize}
    \alert{C的参数都是值传递；传入指针时，复制的是地址值。}
\end{frame}

\begin{frame}{C17中的行为：UB、实现定义与有定义}
    \small 用整数、字符与数组示例，区分语言规则和题目要求。
    \begin{center}
        \footnotesize
        \renewcommand{\arraystretch}{1.2}
        \begin{tabular}{p{0.23\linewidth}p{0.38\linewidth}p{0.27\linewidth}}
            \textbf{类别} & \textbf{标准给了什么保证？} & \textbf{本课例子} \\
            \alert{UB：未定义行为}\newline{\scriptsize Undefined Behavior} & 标准不规定程序应如何表现；不保证报错。 & 有符号溢出、越界访问。 \\
            \alert{实现定义行为}\newline{\scriptsize \mbox{Implementation-defined}\newline Behavior} & 实现作出选择，并在文档中说明。 & 普通\texttt{char}的符号属性。 \\
            \alert{有定义的行为}\newline{\scriptsize Well-defined Behavior} & 按明确的语言规则执行。 & 无符号回绕、合法整数除法。
        \end{tabular}
    \end{center}
    \begin{vspquotered}{有定义，也可能不符合题意}
        \texttt{3 / 2}得到1符合C规则；若题目需要1.5，仍然是逻辑错误。
    \end{vspquotered}
    \footnotesize 这不是穷尽分类：还有\alert{未指定行为}。\alert{编译错误}需先处理，例如使用未声明标识符。
    \note{以C17为准，概念依据C标准术语，不能称为仅有的三类。术语依据N1570第3.4.1、3.4.3、3.4.4节；有定义行为在此作为教学用语，不是穷尽的标准分类。}
\end{frame}

\begin{frame}{有符号溢出与无符号回绕}
    \begin{columns}[T]
        \begin{column}{0.48\textwidth}
            \small \alert{有符号int：UB}
            \lstset{basicstyle=\ttfamily\fontsize{7.6}{9.0}\selectfont,
                numberstyle=\tiny\color{gray},xleftmargin=6pt}
            \begin{Code}*{c}
                \lstinputlisting[lastline=9]{code/behavior/signed-overflow.c}
            \end{Code}
            \small 数学结果是\texttt{INT\_MAX + 1}，超出\texttt{int}范围；不能依赖回绕结果。
        \end{column}
        \begin{column}{0.48\textwidth}
            \small \alert{unsigned int：有定义}
            \lstset{basicstyle=\ttfamily\fontsize{7.6}{9.0}\selectfont,
                numberstyle=\tiny\color{gray},xleftmargin=6pt}
            \begin{Code}*{c}
                \lstinputlisting[lastline=9]{code/behavior/unsigned-wrap.c}
            \end{Code}
            \small 按该无符号类型的取值个数取模，输出\alert{0}；这是语言规定的回绕。
        \end{column}
    \end{columns}
    \medskip
    \small UBSan可报告左侧溢出；右侧通常没有诊断。\alert{改变类型，就可能改变行为分类。}
    \note{依据C17的表达式可表示性规则与无符号运算规则（N1570第6.5第5段、6.2.5第9段）。右侧value += 3在unsigned int域内回绕到0，不属于C标准所指的有符号溢出UB。不要用实跑左侧得到的某个数作为语言保证。}
\end{frame}

\begin{frame}{printf格式必须匹配实参类型}
    \begin{columns}[T]
        \begin{column}{0.48\textwidth}
            \small \alert{格式与实参不匹配：UB}
            \begin{Code}*{c}
                \lstinputlisting[firstline=4,lastline=5,firstnumber=4]{code/behavior/format-mismatch.c}
            \end{Code}
            \small \texttt{value / 4}是\texttt{double}，\texttt{\%d}却要求\texttt{int}；格式串不会替你转换参数。
            \medskip

            修复格式，输出\alert{2.25}：
            \begin{Code}*{c}
                \lstinputlisting[firstline=5,lastline=5,firstnumber=5]{code/behavior/format-fixed.c}
            \end{Code}
        \end{column}
        \begin{column}{0.48\textwidth}
            \small \alert{字符整数提升：有定义}
            \begin{Code}*{c}
                \lstinputlisting[firstline=4,lastline=5,firstnumber=4]{code/behavior/digit-promotion.c}
            \end{Code}
            \small 本课环境中\texttt{char}提升为\texttt{int}，与\texttt{\%c}需要的参数类型匹配。
            \par\medskip
            数字字符连续，\texttt{'3' + 4}对应\texttt{'7'}，因此输出\alert{7}。
        \end{column}
    \end{columns}
    \medskip
    \footnotesize \texttt{-Wall}可提示左侧格式错误，\texttt{-Werror}使本课构建失败；允许编译通过也不消除运行时UB。
    \note{依据N1570第7.21.6.1第9段的格式实参类型匹配规则、6.3.1.1的整数提升及5.2.1的数字字符连续规则。数字字符连续不依赖ASCII，避免把字母连续性误当成C保证。本课演示机是8位char、32位int；常见环境中char提升为int。左侧即使加数值恰好能表示成int，也不会因\%d自动做转换。编译器选择警告或-Werror拒绝构建，与执行该错误调用时的UB规则是两个层面。}
\end{frame}

\begin{frame}{char的符号属性由实现定义}
    \begin{columns}[T]
        \begin{column}{0.46\textwidth}
            \begin{Code}*{c}
                \lstinputlisting[lastline=7]{code/behavior/char-signedness.c}
            \end{Code}
            \small 普通\texttt{char}与\texttt{signed char}、\texttt{unsigned char}是不同的类型。
        \end{column}
        \begin{column}{0.50\textwidth}
            \small 本课8位\texttt{char}环境下，用GCC显式选择：
            \begin{center}
                \renewcommand{\arraystretch}{1.4}
                \begin{tabular}{ll}
                    \textbf{编译选项} & \textbf{输出} \\
                    \texttt{-fsigned-char} & \texttt{-1} \\
                    \texttt{-funsigned-char} & \texttt{255}
                \end{tabular}
            \end{center}
            \begin{itemize}
                \item 默认选择由目标平台与编译器配置决定，并有文档说明。
                \item 这里的差异来自\alert{实现定义的选择}，不是随机结果。
            \end{itemize}
        \end{column}
    \end{columns}
    \medskip
    \small 需要明确的符号属性时，直接使用\texttt{signed char}或\texttt{unsigned char}。
    \note{依据N1570第6.2.5第15段与GCC的C Dialect Options／Integers implementation说明。-1到unsigned char的转换按该无符号类型的模进行，在8位char上得到255，再提升为int后以\%d输出。实际值与位宽有关，不能将255宣传为所有C实现的保证；本页已明确演示环境。选择实现定义行为需看实现文档，单次运行不能证明跨平台结果。}
\end{frame}

\begin{frame}{数组越界与编译错误要分开}
    \begin{columns}[T]
        \begin{column}{0.48\textwidth}
            \small \alert{访问data[6]：UB}
            \begin{Code}*{c}
                \lstinputlisting[firstline=4,lastline=7,firstnumber=4]{code/behavior/bounds.c}
            \end{Code}
            \small 6个元素的合法下标为0至5；\texttt{i == 6}时越界，Sanitizer可帮助定位。
            \par\medskip
            修复循环条件为\texttt{i < 6}，每个元素增加一次，\texttt{data[0]}为1。
        \end{column}
        \begin{column}{0.48\textwidth}
            \small \alert{show中找不到number：编译错误}
            \lstset{basicstyle=\ttfamily\fontsize{7.6}{9.0}\selectfont,
                numberstyle=\tiny\color{gray},xleftmargin=6pt}
            \begin{Code}*{c}
                \lstinputlisting[firstline=3,firstnumber=3,lastline=11]{code/behavior/scope.c}
            \end{Code}
            \small \texttt{main}的局部变量不在\texttt{show}的作用域内；可通过形参传入：\texttt{show(number)}。
        \end{column}
    \end{columns}
    \medskip
    \footnotesize 判断顺序：先检查声明、作用域与类型，再检查运行时边界；\alert{有定义后仍要核对题意。}
    \note{依据N1570第6.5.6的数组访问规则及第6.2.1的作用域规则。右侧使用未声明标识符属于需诊断的翻译错误，不应把它简单归入运行时UB。左侧循环末次访问data[6]是UB，修复后输出1；右侧修复版使用show(int number)与show(number)，输出17。仅运行本次输入没有报告不能证明一般程序没有UB。}
\end{frame}


\begin{frame}{小练习：HW2多项式没有越界，结果仍然错了}
    \begin{columns}[T]
        \begin{column}{0.61\textwidth}
            \begin{Code}*{c}
                \lstinputlisting[firstline=5,lastline=11,firstnumber=5]{code/exercise/polynomial.c}
            \end{Code}
        \end{column}
        \begin{column}{0.35\textwidth}
            仍使用课堂固定数据：\par
            $P(x)=1+2x+3x^2$\par
            $x=2$\par\medskip
            预期\texttt{17.000}，\par
            实际\alert{\texttt{8.000}}。\par\medskip
            这次每个已访问下标都合法，应该检查什么？
        \end{column}
    \end{columns}
    \note{先给学生2至3分钟独立观察。这是根据HW2 P3构造的C语言循环边界练习，不是其原始starter中的错误。}
\end{frame}

\begin{frame}{先写下调试计划，再动手改代码}
    \begin{enumerate}
        \item \alert{复现}：运行\texttt{./build/demo/exercise}，对照预期\texttt{17.000}。
        \item \alert{预测}：下标2、1、0都应被处理，中间结果应依次为3、8、17。
        \item \alert{观察}：在\texttt{evaluate}中跟踪\texttt{i}和\texttt{value}。
        \item \alert{追踪}：执行完下标1后，是否还有下标0这一轮？
        \item \alert{解释}：循环退出条件漏掉了哪个系数？
    \end{enumerate}
    \medskip
    \begin{vspquotered}{给自己一个检查点}
        已经访问的地址都合法，不代表算法完成了所有必要步骤。\par
        修复之后，再用一个能手算的多项式检查你的理解。
    \end{vspquotered}
\end{frame}

\begin{frame}{练习解答：不能漏掉常数项}
    \begin{Code}*{c}
        \lstinputlisting[firstline=5,lastline=11,firstnumber=5]{code/fixed/polynomial.c}
    \end{Code}
    \begin{itemize}
        \item 原来的\texttt{i > 0}跳过了\texttt{coef[0]}对应的最后一步。
        \item 修复为\texttt{i >= 0}，递推结果由8变为$8\times2+1=17$。
        \item 区分两类边界错误：\alert{多访问一步可能越界；少执行一步可能只产生错误结果}。
    \end{itemize}
    \medskip
    修复版\texttt{./build/demo/exercise-fixed}输出\texttt{17.000}。
    \note{C语言案例与练习约25分钟；累加器传递、除法、作业约束与边界练习。最后3页用于5分钟总结与提问。}
\end{frame}

\begin{frame}{下次出错时，照这个顺序开始}
    \begin{enumerate}
        \item 保存源码，记录构建命令，阅读编译警告。
        \item 写出最小失败输入、预期输出与实际输出。
        \item 运行Sanitizer构建；如有报告，定位并解释非法操作。
        \item 对未解释的错误，用GDB观察关键状态，验证具体猜想。
        \item 结合类型、边界、传递方式与生命周期解释原因。
        \item 重新编译，重跑失败输入与边界测试，保留修复记录。
    \end{enumerate}
    \medskip
    \alert{调试记录应该回答：证据是什么？原因是什么？为什么这个修复有效？}
\end{frame}

\begin{frame}{一页速查：先选你需要的观察方式}
    \begin{center}
        \renewcommand{\arraystretch}{1.35}
        \begin{tabular}{ll}
            \textbf{想做什么？} & \textbf{使用什么？} \\
            从固定输入开始 & \texttt{run < input.txt} \\
            在指定位置停止 & \texttt{break file.c:line} \\
            继续／单步／进入／返回 & \texttt{continue / next / step / finish} \\
            查看当前值／自动显示 & \texttt{print / display / info locals} \\
            查看调用关系／选择现场 & \texttt{bt / frame} \\
            只观察某一轮 & \texttt{break file.c:line if condition} \\
            追踪值的变化 & \texttt{watch expression} \\
            检查本课内存与UB案例 & \texttt{-fsanitize=address,undefined}
        \end{tabular}
    \end{center}
    \small GDB内可用\texttt{help 命令名}查帮助；退出使用\texttt{quit}。
\end{frame}

\begin{frame}{课后复现与参考资料}
    \small
    \begin{itemize}
        \item \alert{构建案例}：\texttt{make demos}。
        \item \alert{验证案例}：\texttt{make check-demos}。
        \item \alert{GDB演示脚本}：\texttt{code/commands/session.gdb}。
        \item \alert{修复版与演示说明}：\texttt{code/fixed/}、作业来源说明与本目录\texttt{README.md}。
    \end{itemize}
    \medskip
    \begin{itemize}
        \item \href{https://x-gen-lab.github.io/knowledge-os/02-technologies/embedded/08-toolchain/02-debugging-testing/beginner/02-gdb-basics/}{工程师知识库：GDB调试器基础使用——命令行调试完全指南}。
        \item \href{https://sourceware.org/gdb/current/onlinedocs/gdb.html/}{GDB官方手册}：断点、单步、观察点、调用栈。
        \item \href{https://gcc.gnu.org/onlinedocs/gcc/Instrumentation-Options.html}{GCC官方手册：Instrumentation Options}。
        \item \href{https://www.open-std.org/JTC1/SC22/WG14/www/docs/n1570.pdf}{C标准草案N1570}与\href{https://gcc.gnu.org/onlinedocs/gcc/C-Dialect-Options.html}{GCC字符符号选项说明}。
        \item \href{https://clang.llvm.org/docs/AddressSanitizer.html}{ASan文档}与\href{https://clang.llvm.org/docs/UndefinedBehaviorSanitizer.html}{UBSan文档}。
    \end{itemize}
    \note{最后5分钟总结与提问。PDF内点击参考资料标题即可访问。课件中的报告节选来自配套程序实测；环境差异导致地址、路径与报告格式变化属正常情况。}
\end{frame}
